\section{Conclusion}

L'analyse formelle avec Tamarin a permis d'identifier deux traces d'attaque théoriques sur le
protocole ASKO-OM8464A2. La première concerne l'absence de garantie que l'initiateur $A$ termine
le protocole lorsque le répondant $B$ le termine, et la seconde porte sur la possibilité de
rejouer les messages pour violer la propriété d'accord injectif. Toutefois, ces attaques ne
constituent pas des vulnérabilités dans le cadre de nos spécifications. Le premier cas est
conforme à notre conception où la validation repose sur l'envoi du dernier message par $B$,
et le second cas est exclu par notre règle de vérification des nonces.